% ID: SEL-D0051-OB-E04
% ESTADO: cotejada. Segunda pasada visual y borrador compilado revisados.
% Fuente: Selectividad/sel_2013_matematicas_aplicadas/septiembre/reserva a/EXAMEN 3 -ANDALUCÍA.pdf
% SHA-256: 35f960392892cf4aa4b4a3e5b0a7b6150b323d584bb033f75db6a182b19206d4
% Página: 2; opción B; ejercicio original 4.

\textbf{(2.5 puntos)} Los representantes de un partido político creen que la proporción de sus votantes será al menos del 35\%. Para confirmarlo eligen una muestra al azar de 1200 votantes y obtienen que 336 de ellos son partidarios de votarles. Mediante un contraste de hipótesis, con $H_0:p\geq0.35$, y a un nivel de significación del 0.01, \textquestiondown{}se puede admitir como cierta la creencia de los representantes del partido político?
